\section{At most \texorpdfstring{\(m\)}{m} prime factors}
\label{s-bounded-support}

When \(\omega(n)\leq m\), the support is chosen together with the
exponents. Define
\[
    g_{\integers}^{\leq}(m)
    =
    \sup_{\substack{n\geq5041\\ \omega(n)\leq m}}
    \log G(n).
\]
The CA marginal ordering restricted to the first \(m\) primes solves this
problem, apart from a finite boundary created by \(n\geq5041\).

For \(p\in\{p_1,\ldots,p_m\}\), order all increments \(F(p,k)\),
\(k\geq1\), by decreasing score, processing equal scores as one block.
Starting from \(1\), let \(n_j\) be the integer obtained after the first
\(j\) blocks. Every prefix is a valid exponent vector with at
most \(m\) positive exponents. Direct score comparison gives the first
score-ordered prefixes
\[
    2,\ 6,\ 12,\ 60,\ 120,\ 360,\ 2520,\ 5040,\ 55440,
\]
which agree with Nicolas's CA tables through \(2520\) and at \(5040\) and
\(55440\) \cite[Tables~1--2]{nicolas-2022-sigma}. Thus \(55440\) is the
first score-ordered prefix above Robin's lower cutoff.

\begin{theorem}[Bounded-support characterization]
\label{thm-bounded-support-prefix}
For every \(m\geq6\),
\begin{equation}
\label{e-bounded-support-prefix}
    g_{\integers}^{\leq}(m)
    =
    \max\left\{
        \max_{\substack{5041\leq n<55440\\ \omega(n)\leq m}}
        \log G(n),
        \ 
        \max_{\substack{j\geq1\\ n_j\geq55440}}
        \log G(n_j)
    \right\}.
\end{equation}
Both maxima are attained. Every maximizing integer at least \(55440\)
is a score-ordered prefix.
\end{theorem}

\begin{proof}
We first reduce an arbitrary support to an initial prime segment. Suppose
\[
    n=\prod_{i=1}^{k}q_i^{a_i},
    \qquad
    q_1<\cdots<q_k,
    \qquad
    k\leq m.
\]
Replacing \(q_i\) by \(p_i\) decreases \(n\) and strictly increases
\(\sigma(n)/n\), unless the support was already initial. If the resulting
integer is below \(5041\), multiply it by the smallest power of \(2\)
that reaches \(5041\). The result is below \(10082\), preserves the
support size, and further increases the divisor ratio. For
\(n\geq55440\), this produces a better feasible integer, either in the
finite range in \eqref{e-bounded-support-prefix} or on an initial
support.

An integer on an initial support is a finite selection from the ordered
increment list. Put \(z=\log n\), and choose consecutive prefix integers
with
\[
    \log n_j\leq z\leq\log n_{j+1}.
\]
The fractional-knapsack inequality bounds
\(\log(\sigma(n)/n)\) by the chord joining the corresponding prefix
values. This chord is affine in \(z\), while \(-\log\log z\) is strictly
convex. One of its endpoints therefore has objective at least
\(\log G(n)\). If \(z=\log55440\), then \(n=55440\), which is already a
prefix. If \(z>\log55440\), both adjacent endpoints are feasible. Equality
for a maximizing integer then requires equality in the knapsack bound. If
the selection cuts a tie block, its point lies in the interior of the
affine block segment, where strict convexity makes one endpoint strictly
better.
At an endpoint, equality in the knapsack bound forces every earlier block
to be selected and every later block to be unselected. The total log-size
then forces the intervening block to be either complete or absent. Thus
every maximizing integer in this range is a score-ordered prefix.

The first maximum in \eqref{e-bounded-support-prefix} is finite. Along
the prefixes, the total possible increase in the log divisor ratio is
bounded by
\[
    -\sum_{i=1}^{m}\log\left(1-\frac1{p_i}\right),
\]
whereas \(\log\log\log n_j\) tends to infinity. Hence
\(\log G(n_j)\) tends to minus infinity, and the second maximum is also
attained.
\end{proof}

\begin{corollary}[Unsaturated support]
\label{cor-bounded-support-ca}
Let \(n>55440\) maximize \(g_{\integers}^{\leq}(m)\), and suppose
\(k=\omega(n)<m\). Then \(n\) is the unique colossally abundant number
with parameter
\[
    \eps=\frac1{\log n\log\log n}.
\]
It is also superabundant at its own size, and
\[
    p_k<\log n<p_{k+1}+1,
    \qquad
    k\in\{\pi(\log n)-1,\pi(\log n)\}.
\]
\end{corollary}

\begin{proof}
Let the last selected block have score \(c_-\), and let the next block
have score \(c_+\). On the two adjacent chord intervals, the one-sided
derivatives of the interpolated objective at \(\log n\) are
\[
    c_- -\eps>0,
    \qquad
    c_+ -\eps<0.
\]
The weak inequalities follow from endpoint maximality, and equality is
excluded by strict convexity on the adjacent interval. Thus every
selected score exceeds \(\eps\), while every eligible unselected score is
smaller.

Because \(k<m\), the first-exponent increment for \(p_{k+1}\) is eligible
and unselected.
Monotonicity of \(F\) extends the strict separation to every larger prime
and to all its exponent increments. Equation~\eqref{e-score-dictionary}
then shows that \(n\) uniquely maximizes
\[
    \frac{\sigma(N)}{N^{1+\eps}}
\]
over all positive integers, which is the CA characterization. The same
inequality shows that \(\sigma(N)/N\leq\sigma(n)/n\) whenever \(N\leq n\),
so \(n\) is SA.

The selected first-exponent increment at \(p_k\) and the unselected one
at \(p_{k+1}\) give
\[
    F(p_k,1)>\frac1{\log n\log\log n}>F(p_{k+1},1).
\]
Using
\[
    \frac1{p+1}
    <
    \log\left(1+\frac1p\right)
    <
    \frac1p
\]
and the monotonicity of \(x\log x\) gives the endpoint bounds. The
statement about \(k\) follows.
\end{proof}

The two maxima in \eqref{e-bounded-support-prefix} exchange roles as
\(m\) grows.

\begin{proposition}[Eventual prefix dominance]
\label{prop-eventual-prefix}
For every \(m\geq6\), the first maximum in
\eqref{e-bounded-support-prefix} equals
\(\log G(10080)=-0.0142\ldots\). For \(6\leq m\leq29\), it is the
larger one, \(g_{\integers}^{\leq}(m)=\log G(10080)\), and the unique
maximizer is \(10080=2^{5}3^{2}5\cdot7\). For \(m\geq30\), every
maximizer of \(g_{\integers}^{\leq}(m)\) is a score-ordered prefix
above \(55440\).
\end{proposition}

\begin{proof}
Every integer below \(55440\) has at most six distinct prime factors,
so the first maximum does not depend on \(m\geq6\). A verification over
the window \([5041,55440)\) places it at \(10080\), with the runner-up
\(27720\) lower by \(0.0075\ldots\).

A score-ordered prefix \(n_j\geq55440\) with \(k\) positive exponents
lies in the slice \(\omega(n)=k\), so its objective is at most
\(g_{\integers}(k)\), with equality only at a slice maximizer. For
\(k\leq5\), Proposition~\ref{prop-small-slices} gives
\(g_{\integers}(k)\leq g_{\integers}(4)=\log G(10080)\), with equality
only at \(k=4\), whose maximizer \(10080\) lies below \(55440\), so the
prefix objective is strictly smaller. For \(6\leq k\leq29\), a
verification by the prefix computation of
Theorem~\ref{thm-integer-prefix} gives
\(g_{\integers}(k)<\log G(10080)\). The tightest case is
\(g_{\integers}(29)=-0.0144\ldots\), short by \(0.00012\ldots\). Hence
the second maximum stays strictly below the first for
\(6\leq m\leq29\), which proves the middle claim.

The same computation gives
\(g_{\integers}(30)=-0.0135\ldots>\log G(10080)\). For \(m\geq30\) the
slice \(\omega(n)=30\) is feasible, so
\(g_{\integers}^{\leq}(m)>\log G(10080)\), no integer in the window
attains the maximum, and Theorem~\ref{thm-bounded-support-prefix}
identifies every maximizer as a score-ordered prefix.
\end{proof}

The prefixes must eventually dominate. The slice maxima
\(g_{\integers}(m)\) tend to zero, since the relaxed excess is a
Mertens error by \eqref{e-relaxed-value} and the gap vanishes by
Theorem~\ref{thm-integrality-gap}, while \(\log G(10080)\) is a fixed
negative number. The computation locates the change at \(m=30\).

For a maximizing prefix above \(55440\), the threshold has two
interpretations. If the cap is active, the Lagrangian and size-constrained
statements hold over integers whose prime divisors are at most \(p_m\). If
the cap is inactive, strict separation extends to all primes, so the
optimizer is an unrestricted CA number and hence SA. Thus passing from
\(\omega(n)=m\) to \(\omega(n)\leq m\) restores the first-exponent scores
\(F(p,1)\) as decisions.
